\begin{figure}[hbp]
\centering
\begin{tikzpicture}[
  x=1cm,y=1cm,
  every node/.style={font=\small},
  directed/.style={postaction={decorate},decoration={markings,
    mark=at position 0.56 with {\arrow{Stealth}}}},
  selected/.style={blue!65!black,line width=1pt,directed},
  ordinary/.style={black!60,line width=.7pt,directed}
]
  \node[font=\small\bfseries] at (2,2.0) {Parallel edges of the skeleton};
  \draw[selected] (0,0) .. controls (1,1.5) and (3,1.5) .. (4,0);
  \draw[selected] (0,0) -- (4,0);
  \draw[ordinary] (0,0) .. controls (1,-1.5) and (3,-1.5) .. (4,0);
  \fill (0,0) circle (2pt);
  \fill (4,0) circle (2pt);
  \node[above] at (2,1.13) {$e_j$};
  \node[above] at (2,0) {$e_{j'}$};
  \node at (2,-.60) {$\vdots$};
  \node[below] at (2,-1.13) {other edges};
  \node[left,align=right] at (0,0) {$v(T)=0$};
  \node[right,align=left] at (4,0) {$v(T)=a$};
  \node[align=center] at (2,-2.00)
    {$\ell$ edges, each of length $a$\\cycle $c_j=e_j-e_{j'}$};

  \node[font=\small\bfseries] at (9.4,2.0) {The middle collar on one edge};
  \draw[-{Stealth},black!60] (6.6,1.05) -- (12.2,1.05)
    node[midway,above] {increasing $v(T)$};
  \draw[black!25,line width=2pt] (6.6,0) -- (12.2,0);
  \draw[blue!65!black,line width=2pt] (7.9,0) -- (10.9,0);
  \draw (6.6,-.10) -- (6.6,.10);
  \draw (12.2,-.10) -- (12.2,.10);
  \draw[blue!65!black,fill=white] (7.9,0) circle (2.5pt);
  \draw[blue!65!black,fill=white] (10.9,0) circle (2.5pt);
  \node[below] at (6.6,-.12) {$0$};
  \node[below] at (7.9,-.12) {$D$};
  \node[below] at (10.9,-.12) {$a-D$};
  \node[below] at (12.2,-.12) {$a$};
  \node[above,text=blue!65!black] at (9.4,.12) {$A^\circ$};
  \node[align=center] at (9.4,-1.22)
    {$u_j=c_j'T(1+h_j)$\\$|h_j|<p^{-D}$ on $A^\circ$};
  \node at (9.4,-2.00) {$D>v_p(m)+\dfrac{1}{p-1}$};
\end{tikzpicture}
\caption{The construction in Proposition~\ref{prop:arbitrary-roots},
shown schematically. All edges are oriented by increasing $v(T)$;
the cycle $e_j-e_{j'}$ gives a character unit $u_j$ of slope one
on the annulus corresponding to $e_j$. Division by $m$ on the Jacobian
supplies a root of $u_j$ on the whole inverse annulus. Trimming its
skeletal interval to $D<v(T)<a-D$ makes $h_j$ small enough for
$1+h_j$ to have an $m$-th root, and hence gives a root of $T$.}
\label{fig:parallel-edges}
\end{figure}
